\chapter{FUNDAMENTAÇÃO TEÓRICA}
Este capítulo constitui a base teórica e conceitual do seu TCC. O objetivo é apresentar, analisar e sintetizar as ideias, conceitos, tecnologias e teorias dos autores que fundamentam a sua investigação e sustentam o desenvolvimento da solução proposta.

Diferente de um resumo isolado, a fundamentação teórica deve ser estruturada por afinidade temática, dividindo o texto em seções e subseções que correspondam aos principais conceitos, arquiteturas, métodos ou tecnologias utilizados no trabalho.

Para elaborar a fundamentação teórica é importante:
\begin{uemaball}
    \item abordar os conceitos essenciais e a evolução do tema na literatura científica e técnica;
    \item ater-se às contribuições mais relevantes diretamente relacionadas ao problema e às tecnologias empregadas no projeto;
    \item manter o rigor científico e evitar o plágio acadêmico, referenciando obrigatoriamente toda ideia, definição ou ilustração que não seja de autoria própria.
\end{uemaball}

\section{Seção 1 da fundamentação}
Apresente aqui o primeiro eixo temático da sua fundamentação teórica. Exemplo de citação indireta \cite{towfic2022benchmarking}. Segundo \citeonline{chu2021sustained}, isto é um outro exemplo de citação no texto.

\begin{citacao}
Injeção de dependências (\textit{dependency injection} ou DI) é um padrão de
desenvolvimento de software usado para manter o baixo acoplamento entre classes do
sistema. As dependências de um objeto não são instanciadas programaticamente, mas
sim injetadas de alguma forma \cite{chambers2021macbook}.
\end{citacao}

\subsection{Subseção 1 da fundamentação}
Utilize subseções quando necessário para detalhar tópicos específicos dentro de uma seção da fundamentação teórica, mantendo a organização hierárquica clara do texto.

\section{Trabalhos Correlatos}
Nesta seção devem ser apresentados e analisados trabalhos acadêmicos ou soluções tecnológicas existentes que guardam relação direta com o problema ou com a proposta deste TCC. Recomenda-se descrever o objetivo, a abordagem e os resultados de cada trabalho correlato, destacando ao final uma análise comparativa que evidencie as semelhanças, limitações e os diferenciais da sua proposta em relação às soluções analisadas.
